
\documentclass[11pt]{article}

\usepackage[margin=1in]{geometry}
\usepackage{amsmath,amssymb}
\usepackage{enumitem}
\usepackage{silence}
\usepackage{apacite}
\usepackage{xcolor}
\usepackage[T1]{fontenc}

\title{Review: ``A Unified Blister and Subglacial Hydrology Framework for Supraglacial Lake Drainage Events''}
\author{}
\date{}

\begin{document}

\maketitle

\section{General}

In this paper, the authors analyze the interaction of a supraglacial lake drainage and a
subglacial hydrologic system in a model framework. I really like the approach of this
paper where they address an open question in glaciology using modern fluid mechanics.
The authors bring together the Hewitt (2013) hydrology model with viscous flow under
an elastic sheet (e.g., Hewitt et al., 2015; Tobin and Neufeld, 2023). Overall, I think the
paper is en route to publication after suitable revisions. Below I offer suggestions and
areas where the paper could be improved. One overarching thought is that the paper
could dig deeper: the videos in the supplement are very cool, but I would suggest adding
several more figures to the main text to illustrate the results of the model.

\section{Remarks}

\begin{enumerate}
\item \textcolor{blue}{I am curious about the need for a set prewetting film thickness $h_0$. The blister is
connecting to a hydrologic system outside of the blister zone, so why not make the
prewetting thickness $h$ from the distributed system? I understand that this changes
the dynamics from a constant prewetting film to a variable, but in essence this is
the focus of the paper, correct? I suspect that there are a host of other issues with
this change, and for this reason, it was not included but I think describing these
reasons and discussing the choice is required at a minimum. If including $h$
instead of $h_0$ is possible, that would be ideal.}

Thanks for the suggestion on the pre-wetted layer set-up. As proposed by the reviewer, using the thickness of the cavity sheet, $h$, would resolve the issue of requiring a pre-wetted layer $h_0>0$. However, we think that the cavity-sheet thickness $h$, as an ``effective'' representation of linked cavities, can be a porous sheet that is physically different from the blister sheet that is an actual water sheet. Therefore, we keep the pre-wetted film in our blister model. Future work will be done to find out how the cavity-sheet thickness can be used to regularise the blister front.

\item \textcolor{blue}{The timescales of the blister dynamics and subglacial hydrology appear to be distinct.
The blister evolves as an elastic bending beam whereas the creep closure of
the subglacial system involves viscous creep closure. Some of the simulations are
run out to 10 days and others out to 300 days. A few thoughts:}

\begin{itemize}
\item \textcolor{blue}{A discussion of the Maxwell time for the ice is certainly required. I understand
that this is a first model and elasticity is a more straightforward first step, but
acknowledging the viscoelastic nature of ice is important.}

\item \textcolor{blue}{I wonder if a viscous bending model might be more appropriate and relatively
easy to implement, given the equation correspondence in that regime. This is
the approach of Evatt and Fowler (2007) for work that is similar in spirit to
the current manuscript.}
\end{itemize}

Thanks for suggesting a discussion of Maxwell time. As pointed out by reviewer 1, the characteristic timescale of blister relaxation can be close to the Maxwell time of ice, which ranges from days to weeks. While an elastic bending model can adequately capture the quick formation (hours) of blisters following lake drainage events, the viscous rheology is required to accurately estimate the following relaxation on a weekly timescale. Therefore, a viscoelastic bending model is more appropriate compared with the current elastic model. However, since the model presented here is an initial attempt which leverages many asymptotic results from previous theoretical studies based on the elastic bending theory, we will leave the viscoelastic model for future work.

\item \textcolor{blue}{For the presentation of results, I would rather see $\kappa \neq 0$ through the majority of
the text and then show the sensitivity in figure 7 where $\kappa = 0$ is one of the options.
Otherwise, the logic and flow of the paper are confusing.}

Thanks for the suggestion. We agree that a $\kappa>0$ case would better fit into the logic flow of the manuscript. I've changed figure 7 to the case with $\kappa=10^{-10}$, which is orginally the blue line in figure 8.

\end{enumerate}

\section{Specific Comments}

\begin{enumerate}
\item \textcolor{blue}{Title and key points: these are useful because they convey what was done, but could be more useful if they also conveyed what was learned. I loosely consider the title as a mirror for the abstract and the key points to mirror the conclusions paragraph. For example, the last key point could restate something like the sentence on lines 437--439.}

Thanks for the suggestion. The last key point is changed to "Simulations of a wintertime lake drainage event in western Greenland show consistent trajectories and speed with remote-sensing observations."

\item \textcolor{blue}{Line 197: I would put an ``and'' between 0 and 10.}

This has been added.

\item \textcolor{blue}{Line 207: ``thin cavity''? how thin? quote a number or point to a plot?}

The averaged thickness is approximately $0.03$m. We've added the number to the manuscript.

\item \textcolor{blue}{Line 220: it is interesting that the blister edge is stable and does not result in channels forming along the edge. I wonder if this is because of the choice of $h_0$ rather than $h$ (e.g. Walder, 2010). I also am curious about the time scales for opening versus closing in this model (rapid to open, slow to close).}

Thanks for bringing this up. Yes, the blister edge is stabilised by the pre-wetted layer, as it guarantees finite transmissivity at the blister edge even when the blister thickness is small. In addition, there is a structural reason that channels do not form on the blister edge: channel opening in our model is only driven by dissipation in the cavity sheet and channels, not the blister dissipation. And the blister is coupled with the cavity/channel system only through the leakage term $Q_b$ (eqn. 6). Therefore, the pressure gradient near the blister edge does not directly contribute to channelisation.

In terms of the timescales for opening versus closing, the opening timescale of the blister is controlled by the lake-drainage input, which typically occurs within hours. The closing, on the other hand, is a consequence of both the viscous lateral flow and the leakage, and has been observed to happen on a daily or weekly timescale. Therefore, the model naturally gives ``rapid to open, slow to close'' profile of a blister.

\item \textcolor{blue}{Line 228--230: I think it is worth stating the Stevens et al. (2022) observations here so that I don’t have to go look them up? They are stated later in the paper in more detail but here would be a good place for a preview.}

Thank you for the suggestion. We have added the range of observed front speed (from $10^{-2}~\text{m}~\text{s}^{-1}$ to $1.0~\text{m}~\text{s}^{-1}$) to the end of this sentence.

\item \textcolor{blue}{Figure 2: since it is up-down symmetric, you could consider plotting $h$ on the top half of each panel and $q$ on the bottom half. This would convey a lot of interesting information that isn’t currently in the paper.}

\item \textcolor{blue}{Line 250--251: you should add another figure to show this secondary blister.}

\item \textcolor{blue}{Line 280: I would write this as ``Figures 2 \& 3.''}

The sentence has been modified.

\item \textcolor{blue}{Lines 306--310: why is there an effective viscosity? Is the thought that this value represents a turbulent viscosity? If so, are these reasonable values? This choice deserves more description.}

Thanks for asking about the effective viscosity. Yes, $\mu_{\text{eff}}$ accounts for all characteristics of the ice-bed interface that contribute to the dissipation within the blister, including turbulence. Constraining the value of $\mu_{\text{eff}}$ would require a turbulent model of the blister, which we leave for future study. For the moment, we just pick the value that best matches the observed blister speed.

\item \textcolor{blue}{Lines 325--327: the flow of the results is confusing to me, some leakage seems important and/or inevitable. I think keeping $\kappa \neq 0$ throughout would make the most sense to me.}

Thanks for the helpful suggestion. We agree that allowing some leakage throughout the calculation is more physically consistent, as a blister which is hydrologically isolated from neighbouring cavities and channels is unlikely to exist. We will modify Figure 7 so that it presents a case where the blister loses water while propagating to the ice margin. For Section 4.1, since this subsection is intended to explore the sensitivity of the propagation speed to $\mu_{\text{eff}}$ and $V_b$, we would like to keep $\kappa=0$ to exclude the effect of $\kappa$.

\item \textcolor{blue}{Lines 340--341: would it be possible to make a log-log plot showing the results from your model and comparing them to the outputs from Lai et al. (2021) as well as Tobin and Neufeld (2023).}

Thank you for the suggestion. We agree that a log-log comparison with Lai et al. (2021) and Tobin and Neufeld (2023) would be useful for interpreting blister relaxation. However, a direct quantitative comparison in a single scaling plot is not straightforward because the three models relax through different mechanisms and are controlled by different parameters. Lai et al. (2021) model relaxation through leakage into an underlying porous sheet, parameterised by an effective transmissivity. In contrast, our model focuses on the limiting case of finite leakage into a coupled cavity-channel system that are more complicated. Different from the above, Tobin and Neufeld (2023) consider uplift relaxation by lateral blister propagation, without leakage into surrounding drainage components. So it is inappropriate to cite the model here, as it does not include volume loss due to leakage. We have removed the citation here.

For this reason, we have kept the comparison qualitative in this manuscript, and we have revised the text to clarify which physical mechanism and parameter regime each model represents.

\item \textcolor{blue}{Figure 4: I would change the colour scheme on the right figure to differentiate them from the left. Green means something different on the left and right.}

The colour scheme has been modified.

\item \textcolor{blue}{Lines 385--386: I thought inferring the effective viscosity from the observations was part of the goal, not the reverse?}

Thank you for pointing this out. We agree that the effective viscosity should be constrained from the observed propagation speed. We have revised the wording to clarify this: the observed propagation velocity of approximately \(0.1~\mathrm{m~s^{-1}}\) is used to calibrate \(\mu_{\mathrm{eff}}\), giving a representative value of \(\mu_{\mathrm{eff}}=10~\mathrm{Pa~s}\), rather than prescribing \(\mu_{\mathrm{eff}}\) independently.

\item \textcolor{blue}{Line 401--403: this is a confusing statement in light of the previous sentences.}

Apologies for the confusion here. Since the preceding sentences already discuss the interpretation of the blister-front propagation and the role of the tail region, we have removed this sentence to avoid ambiguity.

\item \textcolor{blue}{There are some capitalization issues in the references, especially the JGR citations.}

Thanks for pointing this issue out. We have 

\end{enumerate}

\section*{References}

\begin{itemize}
\item Evatt, G. W., and A. C. Fowler. \textit{Cauldron subsidence and subglacial floods}. Ann. Glaciol., 45(1):163--168, 2007.

\item Hewitt, I. J. \textit{Seasonal changes in ice sheet motion due to melt water lubrication}. Earth Planet. Sci. Lett., 371:16--25, 2013.

\item Hewitt, I. J., N. J. Balmforth, and J. R. De Bruyn. \textit{Elastic-plated gravity currents}. Eur. J. Appl. Math., 26(1):1--31, 2015.

\item Lai, C.-Y., et al. \textit{Hydraulic transmissivity inferred from ice-sheet relaxation following Greenland supraglacial lake drainages}. Nat. Commun., 12(1):3955, 2021.

\item Stevens, L. A., et al. \textit{Tidewater-glacier response to supraglacial lake drainage}. Nat. Commun., 13(1):6065, 2022.

\item Tobin, S., and J. A. Neufeld. \textit{Evolution of an elastic blister in the presence of sloping topography}. J. Fluid Mech., 967:A5, 2023.

\item Walder, J. S. \textit{R\"othlisberger channel theory: its origins and consequences}. J. Glaciol., 56(200):1079--1086, 2010.
\end{itemize}

\end{document}
